\section{Introduction}

The twin-prime program requires both endpoints to survive every missing prime
filter below a future head. This article investigates a distinct and weaker
program: the first endpoint is square-safe prime, while the shifted endpoint
is required only to avoid primes below a smaller threshold. The conclusion is
a prime plus an integer with at most two prime factors, not a twin-prime pair.

The proof develops the following results in dependency order:

\begin{enumerate}
\item relaxed-weight positivity implies prime-plus-$P_2$ production --- Section~3;
\item the exact divisor local factor and periodic remainder --- Section~4;
\item the exact shifted-divisor prime-progression discrepancy --- Section~5;
\item the exact bilinear character decomposition --- Section~6; and
\item the refuted scalar-density Type-II shortcut --- Section~7.
\end{enumerate}

Sections 8--9 state the remaining analytic program and the exact claim
boundary. No complete-period identity is substituted for an averaged
short-interval distribution theorem.

\subsection{Scope And Evidence Status}

Let $Q$ be a prime head and put

\begin{equation*}
\begin{aligned}
X&=Q^2,
\qquad
W=P(Q)=\prod_{p\lt Q}p.
\end{aligned}
\end{equation*}

The installed-wheel survivors in the square-safe interval are

\begin{equation*}
S_Q
=
\{n\in[Q,Q^2):\gcd(n,W)=1\}.
\end{equation*}

Every $n\in S_Q$ is prime: if it were composite and smaller than $Q^2$, it
would have a prime divisor smaller than $Q$.

We do not prove a new existence theorem here. Classical Chen theory
already supplies infinitely many primes $p$ for which $p+2$ has at most two
prime factors. The project-specific question is whether positivity can be
derived from the Sieve Sequence's own relaxed weights.

We prove the implication in Section~3 and the divisor local factor, cofactor
progression discrepancy, and bilinear character obstruction properties in
Sections~4--6 mathematically; none of the four is yet Stainless-verified. No
statement in this article should be described as formally verified.

The maintained Sieve Sequence construction is Stainless-verified in the
companion chapter article. This article uses that construction as an input but
does not call its new number-theoretic properties verified. Each property
section below states its population and scope before the proof, and links
its supplementary working note at the end.

\subsection{Relation To Known Results}

This article does not reprove Chen's theorem \cite{chen1973}, which already
establishes infinitely many primes $p$ such that $p+2$ has at most two prime
factors. The machinery invoked below is standard in the linear-sieve
literature: the sifted sequence and level-of-distribution framework of
Halberstam and Richert \cite{halberstamrichert1974}, with Type-I/Type-II terminology as in
Iwaniec and Kowalski \cite{iwaniekowalski2004} and Friedlander and Iwaniec \cite{friedlanderiwaniec2010}.
What is project-specific is the population: the square-safe interval
$[Q,Q^2)$ anchored at the head, the nested wheels $2 \mid Z \mid W$, and the
pre-sieved shifted sequence $\mathcal A_Q=\{n+2:n\in S_Q\}$. The modulo-$3$
obstruction proved in Section~7 is a specific fixed-character correlation on the
complete reduced wheel; the classical parity barrier --- a sieve's general
inability to distinguish the parity of the number of prime factors --- is the
broader context for why such local obstructions defeat scalar-density
models, but the two statements are not identical and are not conflated here.
